\documentclass[11pt]{article}
\usepackage{amsmath, amssymb, amsfonts}
\usepackage{bbm}
\usepackage{physics}
\usepackage{geometry}
\usepackage{hyperref}
\geometry{margin=1in}

\title{Exercise Set: Matrices and Applications in Finance}
\author{}
\date{}

\begin{document}
\maketitle

\section{Exercises for Lecture 1: Matrices}

This exercise set follows the structure of the lecture and reinforces the main notions: matrix operations, linear systems, eigenvalues and diagonalization, covariance matrices, portfolio variance, PCA, and matrix exponentials with applications to interest rate dynamics. It is designed to complement the introductory lecture on matrices. The first part (\textbf{Fundamentals}) contains core exercises that we will correct together in a 90-minute session, while the second part (\textbf{Advanced Applications}) is intended for students who wish to explore more sophisticated financial uses of matrix computations, portfolio theory, factor models, and numerical computation with Python (NumPy and PyTorch).


\subsection*{Fundamentals (1h30 in class)}

\paragraph{Exercise 1 — Basic Matrix Operations.}
Let 
\[
A=\begin{pmatrix}2 & -1 \\ 3 & 4\end{pmatrix},\qquad 
B=\begin{pmatrix}1 & 0 \\ 2 & 1\end{pmatrix}.
\]
Compute: (i) $A+B$, (ii) $A-B$, (iii) $AB$, $BA$, (iv) $A^T$, $B^T$, (v) $\det(A)$ and $A^{-1}$.

\vspace{4mm}

Let $v = (3, -1)^T$. Compute $Av$, $Bv$.

\paragraph{Exercise 2 — Image, Kernel and Orthogonality.}
Given
\[
A=\begin{pmatrix}
1 & -1 & 0 \\
-1 & 2 & -1 \\
0 & -1 & 1
\end{pmatrix},
\]
\begin{enumerate}
\item compute the kernel $\ker(A)$ and the image $\mathrm{Im}(A)$;
\item verify the orthogonality relation
\[
\mathrm{Im}(A)= (\ker(A))^{\perp};
\]
\item observe that $A$ is symmetric, that is $A^{\top}=A$;
\item deduce from the fundamental identity
\[
\mathrm{Im}(A^{\top})=(\ker A)^{\perp}
\]
that, for this matrix,
\[
\mathrm{Im}(A)=\mathrm{Im}(A^{\top})=(\ker(A))^{\perp};
\]
\item give a geometric interpretation of your results: which vectors are sent to $0$ by $A$, and on which subspace does $A$ act non-trivially?
\end{enumerate}

\paragraph{Exercise 3 — Bootstrapping Zero-Coupon Prices via a Linear System.}

\textbf{Motivation.}
In fixed-income markets, the most fundamental objects are the \emph{discount factors}
\(
D(T_k)
\),
which represent the present value of receiving \$1 at future maturity \(T_k\).
All yield curve construction procedures---including forward curves and short-rate
models---start from recovering these discount factors from market prices.

A zero-coupon bond of maturity \(T_k\) pays exactly \$1 at time \(T_k\), so by
no-arbitrage its time-0 price \emph{is} the discount factor:
\[
P_k = D(T_k).
\]
Because discounting is cumulative (longer maturities accumulate all earlier
discount effects), it is natural to express discount factors as partial sums of
increments:
\[
D(T_1)=y_1,\qquad D(T_2)=y_1+y_2,\qquad D(T_3)=y_1+y_2+y_3,
\]
which can be written compactly as
\[
D = Ly,
\]
where \(L\) is the unit lower-triangular matrix of ones.  
This representation makes the pricing system \emph{linear and triangular}, the
basis of the standard bootstrapping algorithm used in practice.

\medskip

\textbf{Data.}
You observe the following zero-coupon prices:
\[
P_1=0.97,\qquad P_2=0.92,\qquad P_3=0.855.
\]

\medskip

\textbf{Task.}
\begin{enumerate}
    \item[(1)] Write the pricing relations between \(P_k\) and the increments \(y_1,y_2,y_3\) in matrix form.
    \item[(2)] Identify the resulting linear system \(Ly = P\).
    \item[(3)] Solve the system explicitly for \(y_1,y_2,y_3\) using forward substitution.
\end{enumerate}

\paragraph{Exercise 4 - Finance-Motivated Linear Algebra}
\begin{enumerate}
    \item \textbf{Two-Asset Portfolio Return}\\
    Let asset returns be represented by the vector
    \[
        r = \begin{pmatrix} r_1 \\ r_2 \end{pmatrix},
    \]
    and let the portfolio weights be
    \[
        w = \begin{pmatrix} 0.4 \\ 0.6 \end{pmatrix}.
    \]
    Compute the portfolio return $R_p = w^T r$.   

    \item \textbf{Currency Conversion as a Matrix Transformation}\\
    Suppose you hold two assets: one denominated in USD, one in EUR.  
    The FX conversion matrix to convert values into GBP is
    \[
        C = \begin{pmatrix}
        0.79 & 0 \\
        0 & 0.86
        \end{pmatrix}.
    \]
    Given the holding vector
    \[
        h = \begin{pmatrix} 10{,}000 \\ 8{,}000 \end{pmatrix},
    \]
    compute the value in GBP: $Ch$.

    \item \textbf{Simple Covariance Matrix}\\
    Consider daily centered returns of two assets:
    \[
        X = \begin{pmatrix}
        0.01 & -0.03 \\
        0.02 &  0.01 \\
        -0.01 & 0.00
        \end{pmatrix}.
    \]
    Compute the (unnormalized) covariance matrix $X^T X$.

\end{enumerate}
\paragraph{Exercise 5 — Portfolio Variance.}
Given 
\[
\Sigma = \begin{pmatrix} 0.04 & 0.018 \\ 0.018 & 0.09 \end{pmatrix},\quad 
w=\begin{pmatrix}0.6\\0.4\end{pmatrix},
\]
compute $w^\top \Sigma w$ and interpret contributions of assets and correlation.

\paragraph{Exercise 6 — Diagonalization.}
For 
\[
C=\begin{pmatrix}1 & 0.8\\ 0.8 & 1\end{pmatrix},
\]
compute eigenvalues/eigenvectors.


\paragraph{Exercise 7 — PCA and Dimensionality Reduction (Simplified).}

You are given a dataset of 2D points representing two measured features $(X_1, X_2)$.
The \emph{centered} data matrix is
\[
X =
\begin{pmatrix}
-1 & -1 \\
0 & 0 \\
1 & 1
\end{pmatrix}.
\]
Each row corresponds to one observation $(X_1, X_2)$ after subtracting the sample mean.

\noindent
\textbf{Tasks}
\begin{enumerate}
    \item Compute the empirical covariance matrix of the dataset
    \[
        \widehat\Sigma = \frac{1}{n-1} X^\top X
    \]
    with $n=3$ observations.

    \item Find the eigenvalues and eigenvectors of the covariance matrix.

    \item Identify the principal component (the eigenvector associated with the largest eigenvalue).

    \item Project the data onto the first principal component (reduce from 2D to 1D).  
    Write the resulting 1D coordinates of the three observations.

    \item Briefly explain in words:
    \begin{itemize}
        \item what information is kept when you project onto the first principal component,
        \item what information is lost.
    \end{itemize}

    \item Discuss why PCA is useful here when:
    \begin{itemize}
        \item storage must be reduced (keeping only one coordinate instead of two),
        \item visualization in 1D is needed,
        \item noise is present in directions with very small variance.
    \end{itemize}
\end{enumerate}

\paragraph{Exercise 8 — Empirical PCA with Python}

We observe $n$ realizations of a $d$-dimensional random vector 
$X = (X_1,\dots,X_d)^\top$. The data are stored in a matrix
\[
X = 
\begin{pmatrix}
- & x_1^\top & - \\
- & x_2^\top & - \\
  & \vdots   &   \\
- & x_n^\top & - 
\end{pmatrix}
\in \mathbb{R}^{n \times d},
\]
where each row $x_i^\top$ is one observation.

\medskip

The \emph{empirical mean} of the sample is
\[
\hat \mu = \frac{1}{n} \sum_{i=1}^n x_i \in \mathbb{R}^d,
\]
and the \emph{empirical covariance matrix} is defined by
\[
\widehat{\Sigma} 
= \frac{1}{n-1} \sum_{i=1}^n (x_i - \hat \mu)(x_i - \hat \mu)^\top 
\in \mathbb{R}^{d \times d}.
\]

\medskip

\noindent
\textbf{Setting.}  
You are given an empirical dataset (for example, daily returns of $d$ assets over $n$ days) in a CSV file, or as a numerical table that you can copy into Python.

\medskip

\noindent
\textbf{Tasks}

\begin{enumerate}
    \item Write down explicitly the empirical covariance matrix $\widehat{\Sigma}$ for $d=2$ in terms of the components $X_1$ and $X_2$:
    \[
      \widehat{\Sigma} = 
      \begin{pmatrix}
      \widehat{\mathrm{Var}}(X_1) & \widehat{\mathrm{Cov}}(X_1,X_2) \\
      \widehat{\mathrm{Cov}}(X_2,X_1) & \widehat{\mathrm{Var}}(X_2)
      \end{pmatrix}.
    \]
    Explain how these quantities can be approximated from the sample $\{x_i\}_{i=1}^n$.
    
    \item Using Python and \texttt{NumPy}, write a script that:
    \begin{enumerate}
        \item loads the empirical data into a NumPy array \texttt{X} of shape $(n,d)$ (either from a CSV file using \texttt{np.loadtxt} or \texttt{np.genfromtxt}, or by hard-coding a small matrix);
        \item computes the empirical mean vector $\hat \mu$;
        \item centers the data by subtracting the mean from each observation:
        \[
        X_{\text{centered}} = X - \mathbf{1}_n \hat \mu^\top;
        \]
        \item computes the empirical covariance matrix $\widehat{\Sigma}$ using the formula
        \[
        \widehat{\Sigma} 
        = \frac{1}{n-1} X_{\text{centered}}^\top X_{\text{centered}};
        \]
        \item checks your result by comparing with \texttt{np.cov}.
    \end{enumerate}
    
    \item Still in Python (without using a built-in PCA function):
    \begin{enumerate}
        \item compute the eigenvalues and eigenvectors of $\widehat{\Sigma}$;
        \item sort them in decreasing order of eigenvalues;
        \item identify the first principal component (eigenvector associated with the largest eigenvalue);
        \item project the centered data onto the first one or two principal components, and store the resulting lower-dimensional representation.
    \end{enumerate}
    
    \item Interpret your results:
    \begin{itemize}
        \item What proportion of the total variance is explained by the first principal component?
        \item In your dataset, does reducing the dimension from $d$ to $1$ or $2$ seem reasonable?
        \item In a financial context (e.g.\ asset returns), how can this help with visualization, noise reduction, or factor modeling?
    \end{itemize}
\end{enumerate}


\paragraph{Exercise 9 — Matrix Exponential in a Coupled Linear ODE.}

Consider the linear ODE
\[
\dot{x}(t) = A x(t), 
\qquad
A = \begin{pmatrix}
-0.2 & 0.3 \\
0    & -0.5
\end{pmatrix},
\qquad
x(0) = \begin{pmatrix} 0.1 \\ 0.4 \end{pmatrix}.
\]

\begin{enumerate}
  \item Compute the eigenvalues and eigenvectors of $A$. Is $A$ diagonalizable?
  
  \item Using the diagonalization $A = P D P^{-1}$ (if possible), derive an explicit expression for the matrix exponential
  \[
    e^{At} = P e^{Dt} P^{-1}.
  \]
  Then write $x(t) = e^{At} x(0)$ in closed form.
  
  \item Compare this situation with the case where $A$ is diagonal
  (for instance $A = \mathrm{diag}(-0.2,-0.5)$). 
  Explain why diagonal matrices make the computation of $e^{At}$ particularly simple, 
  and what extra work is needed when $A$ is not diagonal.
  
  \item \textbf{Long-time behavior.}
  \begin{enumerate}
    \item Show that $x(t) \to 0$ as $t \to +\infty$.
    \item Among the two eigenvalues of $A$, identify the one that dominates the long-time dynamics.
    \item Describe qualitatively how the trajectory behaves for large $t$:
    along which direction (which eigenvector) does $x(t)$ decay?
  \end{enumerate}
  
  \item (Optional) Sketch the trajectory of $x(t)$ in the phase plane $(x_1,x_2)$ and indicate the eigenvectors of $A$.
\end{enumerate}


\subsection*{Advanced Applications}

\paragraph{Exercise 10 — Minimum-Variance Portfolio.}
For 
\[
\Sigma=\begin{pmatrix}0.04 & 0.018\\ 0.018 & 0.09\end{pmatrix},
\]
solve 
\[
\Sigma w = \lambda \mathbf{1}
\]
and impose $w_1+w_2=1$.
Interpret the economic meaning of the solution.

\paragraph{Exercise 11 — PCA in a Three-Asset Universe.}
Given 
\[
\Sigma=\begin{pmatrix}
4 & 2 & 0\\
2 & 3 & 0\\
0 & 0 & 1
\end{pmatrix},
\]
compute eigenvalues, identify the dominant factor, and interpret diversification structure.

\paragraph{Exercise 12 — Coupled Vasicek Model and Matrix Differential Equations.}

Consider a two-factor short-rate model for the interest rates
\[
r(t)=
\begin{pmatrix}
r_1(t)\\
r_2(t)
\end{pmatrix},
\]
where the dynamics under the risk–neutral measure are given by the linear ODE system
\[
\frac{d r(t)}{dt}
= 
K\bigl(\theta - r(t)\bigr),
\qquad
K =
\begin{pmatrix}
2 & -1 \\
0 & 3
\end{pmatrix},
\qquad
\theta =
\begin{pmatrix}
0.05\\
0.02
\end{pmatrix}.
\]

The initial value is
\[
r(0)=
\begin{pmatrix}
0.10\\
0.01
\end{pmatrix}.
\]

\bigskip

\textbf{Tasks}

\begin{enumerate}

\item Rewrite the system in the form
\[
\frac{dr(t)}{dt} = A r(t) + b,
\]
and identify explicitly the matrix $A$ and the vector $b$.

\item Show that the solution of this linear ODE system can be written as
\[
r(t)= r(\infty) + e^{At}\bigl(r(0)-r(\infty)\bigr),
\]
where $r(\infty)$ is the long-term mean level of the system.  
Compute $r(\infty)$ explicitly.

\item Compute the matrix exponential $e^{At}$ using the fact that $A$ is triangular.

\item Give the explicit closed form of $r_1(t)$ and $r_2(t)$.

\item Financial interpretation:
\begin{itemize}
\item which factor mean-reverts faster?
\item what is the economic meaning of $r(\infty)$?
\item what happens to $r(t)$ as $t\to\infty$?
\end{itemize}

\item Bond pricing:

The price at time $0$ of a zero-coupon bond with maturity $T$ and instantaneous rate
\[
R(t)=r_1(t)+r_2(t)
\]
is
\[
P(0,T)=
\mathbb{E}\left[
\exp\left(-\int_0^T R(s)\,ds\right)
\right].
\]

Assuming the rates follow the \emph{deterministic} solution you computed above, compute explicitly
\[
\int_0^T R(s)\,ds,
\]
and deduce $P(0,T)$ in closed form.

\end{enumerate}


\paragraph{Exercise 13 - Portfolio Theory and Risk}
\begin{enumerate}
    \item \textbf{Mean–Variance Portfolio}\\
    Let
    \[
        \mu = \begin{pmatrix} 0.05 \\ 0.09 \\ 0.12 \end{pmatrix}, \qquad
        \Sigma =
        \begin{pmatrix}
            0.04 & 0.01 & 0.02 \\
            0.01 & 0.09 & 0.03 \\
            0.02 & 0.03 & 0.16
        \end{pmatrix}.
    \]
    Compute the variance of the portfolio with weights
    \[
        w = (0.3, 0.3, 0.4)^T.
    \]
    (Compute $w^T \Sigma w$.)

    \item \textbf{Uncorrelated Transformation}\\
    Find an orthogonal matrix $Q$ that diagonalizes $\Sigma$, i.e.,
    \[
        Q^T \Sigma Q = D,
    \]
    and interpret this as transforming asset returns into ``uncorrelated'' risk factors.

\end{enumerate}

\paragraph{Exercise 14 - Linear Factor Models}
\begin{enumerate}
    \item Let asset returns follow a one-factor model:
    \[
        r = \alpha + \beta F + \varepsilon.
    \]
    Given a dataset matrix $R$ of returns over $T$ days for $n$ assets (assume $T > n$), express the OLS solution for $(\alpha, \beta)$ in matrix form.

    \item Assume a two-factor model:
    \[
        r_t =
        \begin{pmatrix}
        \beta_{1,1} & \beta_{1,2} \\
        \beta_{2,1} & \beta_{2,2} \\
        \beta_{3,1} & \beta_{3,2}
        \end{pmatrix}
        \begin{pmatrix}
        F_{1,t} \\
        F_{2,t}
        \end{pmatrix}
        + \varepsilon_t.
    \]
    Write the expression for the estimated factors $\hat{F}_t$ that minimize squared reconstruction error (a PCA-like problem).
\end{enumerate}

\paragraph{Exercise 15 — Markov Chains and Probability Viewpoint.}

A company classifies each day the market as being in one of two states:
\[
\text{Bull (B)} \qquad \text{or} \qquad \text{Bear (R)}.
\]

The evolution is modeled by a \emph{Markov chain} with transition matrix
\[
P =
\begin{pmatrix}
0.8 & 0.2 \\
0.3 & 0.7
\end{pmatrix},
\]
where the first row and column correspond to state \textbf{Bull (B)}, and the second row and column correspond to state \textbf{Bear (R)}.

Thus for example
\[
\mathbb{P}(\text{B tomorrow}\mid \text{B today}) = 0.8,
\qquad
\mathbb{P}(\text{R tomorrow}\mid \text{B today}) = 0.2.
\]

Assume that today the market is \textbf{Bull} with probability $1$.

\bigskip

\noindent
\textbf{Tasks}

\begin{enumerate}
    \item Write the Markov property in words for this model, using conditional probabilities
    \[
    \mathbb{P}(X_{n+1} = j \mid X_n = i, X_{n-1},\dots,X_0).
    \]

    \item Compute the probability that the market is \textbf{Bull} tomorrow.

    \item Compute the probability that the market is \textbf{Bull} in two days
    \[
    \mathbb{P}(X_2 = \text{B} \mid X_0=\text{B}).
    \]

    (Hint: use the matrix $P^2$ or condition on the intermediate state $X_1$.)

    \item Find the stationary distribution $\pi=(\pi_B,\pi_R)$ satisfying
    \[
    \pi P = \pi,
    \qquad
    \pi_B+\pi_R=1,
    \qquad
    \pi_B,\pi_R>0.
    \]

    \item Give a probabilistic interpretation of the stationary distribution in one or two sentences.
\end{enumerate}

%%%%%%%%%%%%%%%%%%%%%%%%%%%%%%%%%%%%%%%%%%%%%%%%%%%%%%%%%%%%%%
\end{document}